\section{Portability beyond DMAIC: a two-domain pilot}
\label{app:portability}

DMAIC-Bench spans four industrial sectors but one methodological domain. To test whether the
evaluation design transfers to workflows with different preconditions, we wrote two
life-sciences cases against the unchanged task contract of \S\ref{sec:bench}: the same
guarded decision graph, data steward, separate verifier and scoring axes. Only the domain layer changes: fourteen methods in four new families of the vocabulary, their reference
implementations (numpy and scipy, as for the industrial methods), five public gate predicates,
and task-specific nouns in the instruction (``study team'' for ``plant''). The cases are opt-in at
discovery time and are built as their own dataset, so the 97-task corpus and every number in
the main text are untouched by the pilot.

\paragraph{Case 40: bulk RNA-seq under library-batch structure.}
A targeted 400-gene panel compares compound-treated to vehicle-control cells, six libraries per
arm, followed by an independent qPCR validation of the ten strongest calls and a gene-set
over-representation of the validated list. The scenario reconstructs two published failures:
the mouse ENCODE comparison in which sequencing batch was confounded with
species~\citep{lin2014comparison,gilad2015reanalysis} and the partially nested design of
the ComBat-seq application, in which controls span three batches while each treatment sits in
one (GEO series GSE83083 and GSE59765)~\citep{zhang2020combatseq}.
The brief states four routing preconditions as the facility's rules. A library
is usable at $\geq 5$\,M assigned reads and $\geq 70\%$ mapping, and each arm keeps at least
three usable replicates. These floors are the brief's settings, taken from the low end of
published guidance: three biological replicates as the minimum for inferential analysis, five
million mapped reads as enough for medium to highly expressed genes, and 70 to 90\% as the
expected mapping rate~\citep{conesa2016survey}; the same literature recommends at least six
replicates for whole-transcriptome designs~\citep{schurch2016many}, so three is a floor for
this panel. The treatment contrast must
be estimable separately from batch among the usable libraries. This fails when every batch
holds one arm, because batch and treatment are then confounded~\citep{leek2010tackling}.
With more than one batch, batch enters the model. A candidate list is issued only if the
validation confirms three quarters of the nominated genes; the fraction is the brief's rule, and the qPCR is reported to the MIQE standard~\citep{bustin2009miqe}, which
sets no such threshold.

Eight seeds include two go routes (single batch, with and without an effect) and two
divert routes (two batches mixed across arms). On the divert signal seed, the batch-blind
comparison over-calls by at least ten genes. Two seeds give no-go routes: one arm per batch
or four treated libraries below the read rules. On trap seed~6, the planned sheet is balanced
and the replicate rule passes, but the four failed libraries leave one arm per batch among
the usable ones. On trap seed~7, the panel calls
forty genes convincingly and the qPCR confirms none, so the stop comes after four analyses.
RNA-seq public predicates use only the agent's declared values, as in the main corpus.

\paragraph{Case 41: case-control association with cohort structure.}
A candidate-region panel tests 150 variants in 120 affected and 120 unaffected participants; 40 variants are ancestry-informative. An independent cohort of 60 affected and 60 unaffected participants provides replication on the same panel. The protocol follows the published QC tutorials~\citep{anderson2010data,marees2018tutorial}:
participants with over 5\% missing genotypes are dropped; variants are filtered on call rate, minor-allele frequency and Hardy--Weinberg in
controls; the cohort is checked for ancestry structure on its leading axes before any variant
is tested, and where structure is present but overlapping, the leading principal components
enter as covariates~\citep{price2006principal}; a discovery association is confirmed by
replication~\citep{chanock2007replicating}. The plate rule is the case's own: a plate with more than a tenth of its samples failed is re-genotyped. Eight seeds: two go routes, two divert
routes (mixed ancestry; the unadjusted scan calls at least six variants, the adjusted
scan one or none), two no-go routes (disjoint ancestry on PC1; a failed plate), and two
traps: on seed~6 a balanced recruitment-site split dominates PC1, where status shows no
difference ($p=0.47$), while the status-aligned ancestry split is disjoint on PC2; on seed~7 a
genome-wide-significant discovery hit does not reproduce in replication. Every
fork measures unenforced restraint, as in the four unlocked tasks of
Appendix~\ref{app:steward}: the discovery stages read one batch, and the replication batch remains unlocked.

\paragraph{Construction checks.}
Both cases pass the correspondence checks of Appendix~\ref{app:construction-qc} (C1--C7, no
findings). The reference solution scores reward $1.000$ with perfect routing and numerical
fidelity on every built task it was run on, and the isolation probe finds no truth value in the
agent image.

\paragraph{Contestants and briefs.}
GPT-6-astra, GPT-5.6-sol and GPT-5.6-luna each make three attempts per task using the same scaffold, plain instruction and gateway as in \S\ref{sec:setup}. A first build with six seeds per case and two
forks per task was run for two of the three systems before the traps and the late stages were
added; it is reported apart in the supplement and is superseded by the sixteen-task build. The brief the steward releases first exists in
two forms. The \emph{protocol} brief states each rule and what a failed rule forbids (a design
in which every batch holds one arm supports no contrast; participants of separate ancestry
cannot be compared), which is the regime of the normative instruction of
\S\ref{sec:control}. The \emph{thresholds-only} brief states the same numbers and nothing
about consequences, so the agent must determine whether a confounded design supports a contrast, as it is under the plain instruction in DMAIC; it is run once per task.
Sixteen tasks, two graphs and three systems support per-fork counts, which
Table~\ref{tab:portability} reports with their denominators; they are not population estimates.

\begin{table}[t]
\centering
\small
\caption{\textbf{Cross-domain pilot, full results.} (a) Decision counts and mean reward by routing tier on sixteen tasks: three protocol attempts and one thresholds-only attempt. (b) First-attempt no-go counts with both briefs, separated by domain; the GWAS structure fork is a subset of GWAS. Counts are correct/graded decisions, not independent task samples.}
\label{tab:portability}
\textbf{(a) All recorded attempts}\par\smallskip
\begin{tabular}{llcccccc}
\toprule
system & brief & no-go & divert & go & A-proceed & B-divert & C-stop \\
\midrule
GPT-6-astra & protocol & 22/24 & 12/12 & 102/102 & 0.97 & 0.97 & 0.92 \\
GPT-5.6-sol & protocol & 23/24 & 11/12 & 96/102 & 0.89 & 0.92 & 0.89 \\
GPT-5.6-luna & protocol & 17/24 & 6/12 & 90/102 & 0.88 & 0.78 & 0.57 \\
GPT-6-astra & thresholds & 6/8 & 3/4 & 33/34 & 0.93 & 0.85 & 0.93 \\
GPT-5.6-sol & thresholds & 6/8 & 3/4 & 32/34 & 0.87 & 0.84 & 0.80 \\
GPT-5.6-luna & thresholds & 5/8 & 1/4 & 31/34 & 0.85 & 0.79 & 0.58 \\
\bottomrule
\end{tabular}
\par\medskip
\textbf{(b) Matched first attempts: no-go decisions}\par\smallskip
\begin{tabular}{lcccccc}
\toprule
& \multicolumn{2}{c}{RNA-seq} & \multicolumn{2}{c}{GWAS} & \multicolumn{2}{c}{GWAS structure} \\
\cmidrule(lr){2-3}\cmidrule(lr){4-5}\cmidrule(lr){6-7}
Agent & Protocol & Thresholds & Protocol & Thresholds & Protocol & Thresholds \\
\midrule
GPT-6-astra & 4/4 & 4/4 & 3/4 & 2/4 & 2/2 & 0/2 \\
GPT-5.6-sol & 4/4 & 4/4 & 4/4 & 2/4 & 2/2 & 0/2 \\
GPT-5.6-luna & 4/4 & 4/4 & 1/4 & 1/4 & 0/2 & 0/2 \\
\midrule
Pooled & 12/12 & 12/12 & 8/12 & 5/12 & 4/6 & 0/6 \\
\bottomrule
\end{tabular}
\end{table}


\paragraph{Results.}
Table~\ref{tab:portability} reports counts and denominators. The reference solution solves all sixteen tasks at reward $1.000$. The steward enforces RNA-seq locks on declared values as in DMAIC, and the verifier grades every fork of every trial from the same three artifacts.


Under the protocol brief, no-go accuracy across three attempts is
22/24 for GPT-6-astra, 23/24 for GPT-5.6-sol and 17/24 for GPT-5.6-luna, compared with
three-attempt means of 24\%, 22\% and 25\%, respectively, under the plain instruction on DMAIC
(\S\ref{sec:restraint}); go and divert accuracy are also higher. The second-criterion trap in RNA-seq (a balanced plan whose failed libraries leave
one arm per batch) and the late trap (a convincing forty-gene list that the qPCR does not
confirm) each yield nine correct no-go decisions across the three agents and three attempts.
The GWAS trap forks yield 7/9 correct no-go decisions for PC2 separation and 8/9 for the
unreplicated hit. Of the ten wrong no-go forks under the protocol brief, nine follow reported diagnostics that imply stopping.


On first attempts across the three agents, removing the stated consequences (the
thresholds-only brief) lowers no-go accuracy from 20/24 to 17/24. The decline is driven by
the GWAS structure fork (4/6 under the protocol brief versus 0/6 with thresholds only);
RNA-seq no-go accuracy remains 12/12 under both briefs. Across three protocol attempts,
GPT-5.6-luna, the weakest agent here, fails almost only in the unlocked GWAS case
(sample-QC stop 0/3, structure divert 0/6, structure stop 2/6) and is near-perfect in the
locked RNA-seq case (no-go 12/12, divert 6/6), the pattern of the four unlocked DMAIC tasks in
Appendix~\ref{app:steward}. Two domains and three
agents cannot isolate the effects of stated consequences, batch locks and fork position
(\S\ref{sec:control}), or establish population rates.

